\documentclass[12pt]{article}
\usepackage[ukrainian,shorthands=off]{babel}   % shorthands=off: символ " не активний (потрібно для міток "..." у TikZ всередині \zadtask)
\usepackage{fontspec}
% --- НАЛАШТУВАННЯ ШРИФТІВ ---
% Шрифти Schoolbook-*.otf лежать у корені проєкту. Шлях підбирається автоматично:
% ../ (компіляція з папки теми), ./ (корінь проєкту), ../../ (підпапки теми 28); інакше -- TeX Gyre Schola.
\newcommand{\nmtSchoolbook}[1]{\setmainfont{Schoolbook-Regular.otf}[
    Path = #1,
    BoldFont = Schoolbook-Bold.otf,
    ItalicFont = Schoolbook-Italic.otf,
    BoldItalicFont = Schoolbook-BoldItalic.otf
]}
\IfFileExists{../Schoolbook-Regular.otf}{\nmtSchoolbook{../}}{%
 \IfFileExists{./Schoolbook-Regular.otf}{\nmtSchoolbook{./}}{%
  \IfFileExists{../../Schoolbook-Regular.otf}{\nmtSchoolbook{../../}}{%
   \setmainfont{texgyreschola}[Extension=.otf, UprightFont=*-regular, BoldFont=*-bold, ItalicFont=*-italic, BoldItalicFont=*-bolditalic]}}}
\usepackage[italic]{mathastext}
\usepackage[a4paper,margin=1.8cm,bottom=2cm,top=2cm]{geometry}
\usepackage{amsmath,amssymb}
\usepackage{tikz}
\usepackage{pgfplots}
\pgfplotsset{compat=1.18}
\usepgfplotslibrary{fillbetween}
\usetikzlibrary{calc,patterns,angles,quotes,intersections,babel,3d,shapes.symbols,shapes.geometric,decorations.pathreplacing,shadings,arrows.meta,hobby}
\usepackage{xcolor,array,fancyhdr,enumitem,multicol}
\usepackage{colortbl,multirow,diagbox,needspace}
\usepackage{tcolorbox}
\tcbuselibrary{skins,breakable}
\usepackage[hidelinks]{hyperref}
% водяний знак; щоб зібрати без нього (версія для вчителів), перед \documentclass пишуть \def\nmtnowatermark{}
\ifdefined\nmtnowatermark\else
\usepackage[firstpage=false, text=@pvtr2525, scale=0.6, color=gray!12]{draftwatermark}
\fi

% --- АВТОСТИСКАННЯ: рисунок/сітка, ширша за колонку, масштабується до ширини колонки ---
\newsavebox{\nmtfitbox}
\newenvironment{nmtfit}{\begin{lrbox}{\nmtfitbox}}{\end{lrbox}%
  \ifdim\wd\nmtfitbox>\linewidth \resizebox{\linewidth}{!}{\usebox{\nmtfitbox}}\else\usebox{\nmtfitbox}\fi}
\let\nmtIncludegraphicsOrig\includegraphics
\renewcommand{\includegraphics}[2][]{\begin{nmtfit}\nmtIncludegraphicsOrig[#1]{#2}\end{nmtfit}}


\definecolor{mainGreen}{RGB}{34, 120, 64}
\definecolor{yearOrange}{RGB}{220, 100, 30}
\definecolor{yearcolor}{RGB}{220, 100, 30}
\definecolor{titleBg}{RGB}{225, 240, 220}
\definecolor{instrBg}{RGB}{255, 240, 220}
\definecolor{instrBorder}{RGB}{220, 150, 80}
\definecolor{headerblue}{RGB}{0, 102, 204}   % використовується в рисунках старої бази

\setlength{\parindent}{0pt}
\emergencystretch=1.5em

\pagestyle{fancy}
\fancyhf{}
\renewcommand{\headrulewidth}{0.4pt}
\renewcommand{\footrulewidth}{0.4pt}
\fancyhead[C]{\small\color{gray!80} Пробний НМТ № 1 (повний)}
\fancyhead[R]{\small\color{mainGreen}\textbf{\href{https://t.me/pvtr2525}{@pvtr2525}}}
\fancyfoot[L]{\small\color{gray!80}\href{https://t.me/pvtr2525}{ТГ: @pvtr2525}}
\fancyfoot[C]{\small\color{gray!80}\thepage}
\fancyfoot[R]{\small\color{gray!80}НМТ з математики}

% --- ТАБЛИЦІ ВІДПОВІДЕЙ ---
% ширина клітинки = (ширина рядка - відступи) / 5: на всю сторінку це ~3 см, у minipage поруч із рисунком таблиця стискається;
% п'ять колонок |m{w}| мають 10 відступів \tabcolsep і 6 вертикальних ліній -- тоді таблиця рівно на \linewidth
\newlength{\nmtcellw}
\newcommand{\nmtSetCell}{\setlength{\nmtcellw}{\dimexpr(\linewidth-10\tabcolsep-6\arrayrulewidth)/5\relax}}
\newcommand{\answerTable}[5]{%
\par\nopagebreak\vspace{0.15cm}
\noindent\nmtSetCell
\begin{tabular}{|*{5}{>{\centering\arraybackslash}m{\nmtcellw}|}}
\hline
\rule[-0.15cm]{0pt}{0.6cm}\textbf{А} & \rule[-0.15cm]{0pt}{0.6cm}\textbf{Б} & \rule[-0.15cm]{0pt}{0.6cm}\textbf{В} & \rule[-0.15cm]{0pt}{0.6cm}\textbf{Г} & \rule[-0.15cm]{0pt}{0.6cm}\textbf{Д} \\
\hline
\rule[-0.2cm]{0pt}{0.75cm}#1 & \rule[-0.2cm]{0pt}{0.75cm}#2 & \rule[-0.2cm]{0pt}{0.75cm}#3 & \rule[-0.2cm]{0pt}{0.75cm}#4 & \rule[-0.2cm]{0pt}{0.75cm}#5 \\
\hline
\end{tabular}
\par\vspace{0.25cm}
}
\newcommand{\answerTableTall}[5]{%
\par\nopagebreak\vspace{0.15cm}
\noindent\nmtSetCell
\begin{tabular}{|*{5}{>{\centering\arraybackslash}m{\nmtcellw}|}}
\hline
\rule[-0.2cm]{0pt}{0.7cm}\textbf{А} & \rule[-0.2cm]{0pt}{0.7cm}\textbf{Б} & \rule[-0.2cm]{0pt}{0.7cm}\textbf{В} & \rule[-0.2cm]{0pt}{0.7cm}\textbf{Г} & \rule[-0.2cm]{0pt}{0.7cm}\textbf{Д} \\
\hline
\rule[-0.45cm]{0pt}{1.2cm}#1 & \rule[-0.45cm]{0pt}{1.2cm}#2 & \rule[-0.45cm]{0pt}{1.2cm}#3 & \rule[-0.45cm]{0pt}{1.2cm}#4 & \rule[-0.45cm]{0pt}{1.2cm}#5 \\
\hline
\end{tabular}
\par\vspace{0.25cm}
}
\newcommand{\instructionBox}[1]{%
\par\vspace{0.3cm}
\begin{tcolorbox}[colback=instrBg, colframe=instrBorder, boxrule=0.6pt, arc=2pt,
    left=8pt, right=8pt, top=4pt, bottom=4pt]
\centering\small\bfseries #1
\end{tcolorbox}
\par\vspace{0.15cm}
}
\newcommand{\sectionTitle}[1]{%
\par\vspace{0.4cm}
\begin{tcolorbox}[colback=titleBg, colframe=titleBg, boxrule=0pt, arc=8pt,
    halign=center, fontupper=\Large\bfseries\color{mainGreen}]
\hyphenpenalty=10000 \exhyphenpenalty=10000 #1
\end{tcolorbox}
\par\vspace{0.3cm}
}
\newcommand{\task}[2]{\noindent\textbf{#1.}\ \ #2\par\vspace{0.2cm}}
% --- НМТ-СТИЛЬ ПОЛЕ ВІДПОВІДІ ---
\newcommand{\nmtAnswerBox}{%
\par\nopagebreak\vspace{0.2cm}\noindent
Відповідь:\ \
\begin{tabular}{@{}*{4}{|>{\centering\arraybackslash}p{0.8cm}}|@{\hspace{0.1cm}\textbf{,}\hspace{0.1cm}}*{3}{|>{\centering\arraybackslash}p{0.8cm}}|@{}}
\hline
\rule[-0.2cm]{0pt}{0.7cm} & & & & & & \\
\hline
\end{tabular}
\par\vspace{0.3cm}
}
% --- СІТКА ДЛЯ МАТЧИНГУ ---
\newsavebox{\nmtgridbox}
\newcommand{\matchingGrid}{%
    \sbox{\nmtgridbox}{\matchingGridRaw}%
    \ifdim\wd\nmtgridbox>\linewidth \resizebox{\linewidth}{!}{\usebox{\nmtgridbox}}\else\usebox{\nmtgridbox}\fi}
\newcommand{\matchingGridRaw}{%
    \begingroup
    \renewcommand{\arraystretch}{1.15}
    \setlength{\tabcolsep}{4pt}
    \footnotesize
    \begin{tabular}{r|c|c|c|c|c|}
         \multicolumn{1}{c}{} & \multicolumn{1}{c}{\textbf{А}} & \multicolumn{1}{c}{\textbf{Б}} & \multicolumn{1}{c}{\textbf{В}} & \multicolumn{1}{c}{\textbf{Г}} & \multicolumn{1}{c}{\textbf{Д}} \\ \cline{2-6}
         \textbf{1} & \phantom{X} & \phantom{X} & \phantom{X} & \phantom{X} & \phantom{X} \\ \cline{2-6}
         \textbf{2} & & & & & \\ \cline{2-6}
         \textbf{3} & & & & & \\ \cline{2-6}
    \end{tabular}
    \endgroup
}
% --- ДОДАТКОВІ МАКРОСИ ЗБІРНИКА ---
\newcommand{\taskBlock}[1]{%
\par\vspace{0.45cm}
\noindent{\color{mainGreen}\rule{\linewidth}{0.8pt}}\par\vspace{0.12cm}
\noindent{\small\bfseries\color{mainGreen}#1}\par\vspace{0.25cm}
}
\newcounter{zad}
\newcommand{\zadtask}[1]{\stepcounter{zad}\noindent\textbf{\thezad.}\ \ #1\par\nopagebreak\vspace{0.2cm}}
\newcommand{\zadnum}{\stepcounter{zad}\textbf{\thezad.}\ \ }
\newcounter{ans}
\newcommand{\ansitem}[1]{\stepcounter{ans}\noindent\textbf{\theans.}\ #1\par\vspace{0.07cm}}
% мітка року: нерозривна, притиснута праворуч; якщо не вміщається -- праворуч на новому рядку
\newcommand{\nmtyear}[1]{\unskip\nobreak\finalhyphendemerits=0 \hfill\penalty100\hskip1em\hbox{}\nobreak\hfill{\small\color{yearOrange}\mbox{\rule{0pt}{2.3ex}(НМТ~#1)}}}
% --- МАКРОС КОРОТКОГО РОЗВ'ЯЗКУ ---
\newcommand{\solution}[1]{%
\par\vspace{0.12cm}
\begin{tcolorbox}[colback=titleBg!45, colframe=mainGreen!55, boxrule=0.4pt, arc=2pt,
    left=6pt, right=6pt, top=3pt, bottom=3pt, breakable]
\footnotesize\textbf{Короткий розв'язок.}\ #1
\end{tcolorbox}
\par\vspace{0.12cm}
}
% Розв'язки й відповіді (є до завдань НМТ-2026): \showsolutionstrue -- показати, \showsolutionsfalse -- сховати
\newif\ifshowsolutions
\showsolutionsfalse
% --- ЗАГОЛОВКИ З ПУНКТАМИ КЛІКАБЕЛЬНОГО ЗМІСТУ ---
\setcounter{tocdepth}{2}
\newcommand{\chapterTitle}[1]{%
\clearpage\phantomsection\addcontentsline{toc}{section}{#1}%
\sectionTitle{#1}}
\newcommand{\typeTitle}[1]{%
\needspace{6\baselineskip}%
\phantomsection\addcontentsline{toc}{subsection}{\quad #1}%
\par\vspace{0.3cm}
\noindent{\large\bfseries\color{yearOrange}#1}\par\vspace{0.08cm}
\noindent{\color{yearOrange}\rule{\linewidth}{0.4pt}}\par\nopagebreak\vspace{0.2cm}}
\raggedbottom
% усередині одного завдання сторінка не розривається: нерозривний вертикальний проміжок
% і нерозривна межа абзаців (умова ніколи не відділяється від варіантів, рисунка чи сітки)
\newcommand{\nbvspace}[1]{\nopagebreak\vspace{#1}}
\newcommand{\nmtnobreak}{\ifvmode\penalty10000 \fi}
% --- ЄДИНИЙ МАКЕТ ЗАВДАНЬ НА ВІДПОВІДНІСТЬ ---
% мітка в колонці сталої ширини, текст -- у parbox: продовження довгого рядка
% вирівнюється під текстом, а не під міткою; відступ між пунктами однаковий скрізь
\newlength{\nmtMatchLab}\setlength{\nmtMatchLab}{1.6em}
\newlength{\nmtMatchSep}\setlength{\nmtMatchSep}{0.28cm}
\newcommand{\matchHead}[1]{\par\noindent\textit{#1}\par\nopagebreak\vspace{0.2cm}}
% шаблонний заголовок стовпця зі старого макета -- лише коли завдання не має власного \matchHead
\makeatletter\newcommand{\nmtHeadUnless}[2]{\in@{\matchHead}{#1}\ifin@\else#2\fi}\makeatother
\newcommand{\matchItem}[2]{\par\noindent\makebox[\nmtMatchLab][l]{\textbf{#1}}%
\parbox[t]{\dimexpr\linewidth-\nmtMatchLab\relax}{\raggedright #2}\par\nopagebreak\vspace{\nmtMatchSep}}
\newcommand{\ansTheme}[1]{\par\vspace{0.25cm}\noindent{\bfseries\color{mainGreen}#1}\par\vspace{0.12cm}}
\newcommand{\ansType}[1]{\par\vspace{0.1cm}\noindent{\itshape\color{yearOrange}#1}\par\vspace{0.08cm}}

\graphicspath{{../1. Числа. Звичайні та десяткові дроби. Модуль числа/}{../10. Основи планіметрії/}{../11. Трикутники/}{../12. Паралелограм/}{../13. Прямокутник/}{../14. Квадрат/}{../15. Ромб/}{../16. Трапеція/}{../17. Коло, круг і їх елементи/}{../18. Координати і вектори на площині і в просторі. Рівняння кола/}{../19. Лінійні та квадратні нерівності. Системи лінійних нерівностей. Метод інтервалів/}{../2.відсотки, арифметичні задачі, подільність/}{../20. Функція та її властивості/}{../21.  Графіки елементарних функцій, їх побудова графіків функцій та їх перетворення/}{../22. Модульні  рівняння та нерівності/}{../23. Ірраціональні рівняння та нерівності/}{../24. Арифметична прогресія/}{../25. Геометрична прогресія/}{../26. Тригонометричні вирази/}{../27. Тригонометричні рівняння/}{../28. Показникова функція. Показникові рівняння/Показникові рівняння/}{../28. Показникова функція. Показникові рівняння/Показникові функція і вирази/}{../29. Показникові нерівності/}{../3. Степені. Одночлени. стандартний вигляд числа/}{../30. Логарифм. Логарифмічна функція/}{../31. Логарифмічні рівняння/}{../32. Логарифмічні нерівності/}{../33. Похідна/}{../34. Первісна/}{../35. Комбінаторика/}{../36. Теорія ймовірності/}{../37. Статистика/}{../38. Призма. Паралелепіпед. Куб/}{../39. Піраміда/}{../4.Многочлени. ФСМ/}{../40. Циліндр/}{../41. Конус/}{../42. Куля. Сфера/}{../43. Параметри/}{../5. дробові раціональні вирази і перетворення/}{../6. Лінійні рівняння/}{../7. Квадратні корені. Корені вищих степенів. Раціональний степінь/}{../8. квадратні рівняння, і рівняння, що до них зводяться/}{../9. Системи рівнянь/}}
\renewcommand{\instructionBox}[1]{%
\par\vspace{0.3cm}\noindent
\begin{tcolorbox}[nobeforeafter, width=\linewidth, colback=instrBg, colframe=instrBorder, boxrule=0.6pt, arc=2pt,
    left=8pt, right=8pt, top=4pt, bottom=4pt]
\centering\small\bfseries #1
\end{tcolorbox}
\par\nopagebreak\vspace{0.15cm}\nopagebreak
}
\renewcommand{\nmtAnswerBox}{%
\par\nopagebreak\vspace{0.2cm}\noindent
Відповідь:\ \
\begin{tabular}{@{}*{5}{|>{\centering\arraybackslash}p{0.8cm}}|@{\hspace{0.1cm}\textbf{,}\hspace{0.1cm}}*{3}{|>{\centering\arraybackslash}p{0.8cm}}|@{}}
\hline
\rule[-0.2cm]{0pt}{0.7cm} & & & & & & & \\
\hline
\end{tabular}
\par\vspace{0.3cm}
}

\begin{document}
\chapterTitle{Пробний НМТ з математики № 1 (повний): відповіді}

\begin{center}\renewcommand{\arraystretch}{1.25}
\begin{tabular}{|c|c||c|c||c|c|}\hline
\textbf{№} & \textbf{Відповідь} & \textbf{№} & \textbf{Відповідь} & \textbf{№} & \textbf{Відповідь} \\ \hline

1 & А & 9 & Д & 17 & 1~--~В, 2~--~Б, 3~--~Г \\ \hline
2 & А & 10 & В & 18 & 1~--~В, 2~--~Г, 3~--~Б \\ \hline
3 & Б & 11 & Б & 19 & $32$ \\ \hline
4 & Б & 12 & Г & 20 & $96$ \\ \hline
5 & Д & 13 & Г & 21 & $500$ \\ \hline
6 & Г & 14 & В & 22 & $2$ \\ \hline
7 & В & 15 & Д &  &  \\ \hline
8 & А & 16 & 1~--~Д, 2~--~В, 3~--~А &  &  \\ \hline
\end{tabular}\end{center}

\noindent{\small\color{gray!80!black}Оцінювання: 1--15 -- по 1 балу; 16--18 -- по 1 балу за кожну правильно встановлену відповідність (до 3 балів); 19--22 -- по 2 бали. Разом 32 бали.}\par\vspace{0.3cm}

\sectionTitle{Розв'язки й типові помилки}

\par\noindent\textbf{1.} $60\,\%$ -- за $36$~хв, тобто $1\,\%$ -- за $0{,}6$~хв; $20\,\%$ -- за $12$~хв.\quad{\small\color{gray!80!black}(Рух і продуктивність)}\par\nopagebreak

\noindent{\footnotesize\color{gray!85!black}Типові помилки: Б -- вважали, що останні $20\,\%$ заряджаються стільки ж, скільки попередні $60\,\%$; В -- $60\,\%$ -- це $6$ частин по $10\,\%$, а $20\,\%$ вважали однією частиною; Г -- склали пропорцію з $80\,\%$: $36\cdot\frac{20}{80}$; Д -- знайшли час від $20\,\%$ до $100\,\%$.}\par\vspace{0.15cm}

\par\noindent\textbf{2.} Щонайменше $10$~тисяч -- у вівторок, середу, п'ятницю й суботу: $4$ дні.\quad{\small\color{gray!80!black}(Стовпчикова діаграма)}\par\nopagebreak

\noindent{\footnotesize\color{gray!85!black}Типові помилки: Б -- зарахували день із $9$~тисячами кроків (майже ціль); В -- пропустили один із двох днів рівно з $10$~тисячами; Г -- рахували лише дні, коли кроків було \textit{більше} за $10$~тисяч; Д -- рахували дні щонайменше з $8$~тисячами кроків.}\par\vspace{0.15cm}

\par\noindent\textbf{3.} $O$ -- середина дошки, тому $OA_1=OB_2$ і трикутник $A_1OB_2$ рівнобедрений: $\angle OB_2A_1=\angle OA_1B_2=24^\circ$. Точки $A_1$, $O$, $B_1$ лежать на одній прямій, тож $\alpha=\angle B_1OB_2$ -- зовнішній кут трикутника $A_1OB_2$: $\alpha=24^\circ+24^\circ=48^\circ$. (Інакше: у новому положенні дошка теж утворює із землею кут $24^\circ$, тож вона повернулася на $24^\circ+24^\circ=48^\circ$.)\quad{\small\color{gray!80!black}(Практична задача: гойдалка-балансир (кут повороту, рівнобедрений трикутник))}\par\nopagebreak

\noindent{\footnotesize\color{gray!85!black}Типові помилки: А -- вважали, що дошка повернулася на кут, який вона утворює із землею; В -- знайшли кут між дошкою і вертикаллю: $90^\circ-24^\circ$; Г -- знайшли кут $A_1OB_2$ між опущеними половинами дошки: $180^\circ-2\cdot24^\circ$; Д -- узяли кут, суміжний з кутом $24^\circ$: $180^\circ-24^\circ$.}\par\vspace{0.15cm}

\par\noindent\textbf{4.} $\lg1=0$: $-3(x-2)>6$, $x-2<-2$, $x<0$; найбільший цілий розв'язок $-1$.\quad{\small\color{gray!80!black}(Лінійна нерівність: найменший цілий розв'язок чи кількість цілих)}\par\nopagebreak

\noindent{\footnotesize\color{gray!85!black}Типові помилки: А -- вважали $\lg1=1$: $x<-1$; В -- включили межу $x=0$, хоча нерівність строга; Г -- не змінили знак при діленні на $-3$ і взяли найменший цілий; Д -- не змінили знак: $x>0$, найбільшого немає.}\par\vspace{0.15cm}

\par\noindent\textbf{5.} Площини $SBC$ і $ABC$ перетинаються по прямій $BC$. $SM\perp BC$ (медіана рівнобедреного трикутника $SBC$ є висотою), $OM\parallel AB$ (середня лінія трикутника $ABC$), а $AB\perp BC$, тому й $OM\perp BC$. Кут між площинами -- це кут між перпендикулярами до прямої їх перетину, проведеними в цих площинах через одну точку: $\angle SMO$.\quad{\small\color{gray!80!black}(Кути в просторі: практичні задачі)}\par\nopagebreak

\noindent{\footnotesize\color{gray!85!black}Типові помилки: А -- узяли кут між бічним \textit{ребром} $SC$ і площиною основи, а не між гранню й основою; Б -- узяли кут між бічним ребром і стороною основи -- кут самої грані $SBC$; В -- узяли кут між апофемою $SM$ і прямою перетину площин $BC$ (він прямий), а треба -- між двома перпендикулярами до $BC$; Г -- узяли кут між апофемою й висотою піраміди, а не між апофемою й площиною основи.}\par\vspace{0.15cm}

\par\noindent\textbf{6.} $9x^2-y^2=(3x-y)(3x+y)$, тому (оскільки $3x+y\neq0$) вираз дорівнює $\dfrac{1}{3x-y}$. З $6x-2y=5$: $3x-y=\dfrac52$, отже, значення $\dfrac25$.\quad{\small\color{gray!80!black}(Значення виразу за відомим співвідношенням)}\par\nopagebreak

\noindent{\footnotesize\color{gray!85!black}Типові помилки: А -- розклали знаменник з помилкою знака: $9x^2-y^2=(y-3x)(y+3x)$; Б -- не поділили на $2$ і не перевернули дріб: відповіддю взяли дане число $5$; В -- вважали, що $3x-y=5$ (не поділили $6x-2y=5$ на $2$); Д -- після скорочення записали $3x-y=\frac52$ замість $\frac{1}{3x-y}$.}\par\vspace{0.15cm}

\par\noindent\textbf{7.} $0{,}16=\frac{4}{25}=\left(\frac25\right)^2$, $\frac{125}{8}=\left(\frac52\right)^3=\left(\frac25\right)^{-3}$: $\left(\frac25\right)^{2x}=\left(\frac25\right)^{-3}$, $2x=-3$, $x=-1{,}5\in[-2;\,-1)$.\quad{\small\color{gray!80!black}(Показникове: зведення до однієї основи (дробова чи десяткова основа))}\par\nopagebreak

\noindent{\footnotesize\color{gray!85!black}Типові помилки: А -- степінь степеня: $\left(\left(\frac25\right)^2\right)^x=\left(\frac25\right)^{2+x}$ -- додали показники, $x=-5$; Б -- загубили квадрат: $0{,}16^x=\left(\frac25\right)^{x}$, тоді $x=-3$; Г -- з $2x=-3$ поділили навпаки: $x=-\frac23$; Д -- вважали $\frac{125}{8}=\left(\frac25\right)^3$ (не перевернули дріб): $2x=3$, $x=1{,}5$.}\par\vspace{0.15cm}

\par\noindent\textbf{8.} Спуск змінюється підйомом при $x=10$ і $x=30$.\quad{\small\color{gray!80!black}(Точки екстремуму за графіком)}\par\nopagebreak

\noindent{\footnotesize\color{gray!85!black}Типові помилки: Б -- точки максимуму; В -- лише найнижча точка гірки; Г -- висоти в точках мінімуму, а не значення $x$; Д -- кінець проміжку не є точкою екстремуму.}\par\vspace{0.15cm}

\par\noindent\textbf{9.} I і II правильні. III хибне: навпіл діляться лише діаметри.\quad{\small\color{gray!80!black}(Коло: хорди, діаметри, дотичні, кути)}\par\nopagebreak

\noindent{\footnotesize\color{gray!85!black}Типові помилки: А -- не знали зв'язку центрального й вписаного кутів; Б -- вважали, що хорди діляться навпіл; В -- сплутали властивість хорди з властивістю діаметрів; Г -- властивість діаметрів поширили на будь-які хорди.}\par\vspace{0.15cm}

\par\noindent\textbf{10.} $a_7-a_3=4d=-12$, тож $d=-3$. $a_n=26-3(n-1)=29-3n>0$ при $n<9\frac23$, тобто $n\le9$: додатних членів $9$ ($a_9=2$, $a_{10}=-1$).\quad{\small\color{gray!80!black}(Арифметична прогресія: член за двома відомими)}\par\nopagebreak

\noindent{\footnotesize\color{gray!85!black}Типові помилки: А -- не поділили на кількість кроків: $d=-12$ (члени $26;\ 14;\ 2;\ -10;\ \ldots$); Б -- записали $a_n=a_1+nd$: $26-3n>0$, $n<8\frac23$; Г -- з $n<9\frac23$ узяли $n=10$ -- це номер першого від'ємного члена; Д -- з $a_3=a_7+12$ узяли $d=3$ (знак різниці): прогресія зростає.}\par\vspace{0.15cm}

\par\noindent\textbf{11.} Непарних цифр п'ять: $1$, $3$, $5$, $7$, $9$. Для першого коліщатка -- $5$ варіантів, для другого -- $4$ (цифри різні), далі $3$ і $2$: $A_5^4=5\cdot4\cdot3\cdot2=120$ кодів.\quad{\small\color{gray!80!black}(Розміщення: вираз для кількості)}\par\nopagebreak

\noindent{\footnotesize\color{gray!85!black}Типові помилки: А -- не врахували порядок цифр -- так рахують вибір набору цифр, а не код; В -- не врахували, що цифри непарні: вибирали з усіх $10$ цифр; Г -- не врахували, що цифри різні: на кожне коліщатко брали будь-яку з $5$ цифр; Д -- додали кількості варіантів для кожного коліщатка замість перемножити.}\par\vspace{0.15cm}

\par\noindent\textbf{12.} Радіус -- відстань від центра до площини $xy$, тобто $|z|=5$; діаметр $10$.\quad{\small\color{gray!80!black}(Відстань між точками; сфера з відомим центром)}\par\nopagebreak

\noindent{\footnotesize\color{gray!85!black}Типові помилки: А -- знайшли радіус, а не діаметр; Б -- узяли відстань до площини $yz$: $2|x|$; В -- узяли відстань до площини $xz$: $2|y|$; Д -- подвоїли відстань від центра до початку координат.}\par\vspace{0.15cm}

\par\noindent\textbf{13.} Для тупого кута $\sin\alpha>0$, а $\cos\alpha<0$ і $\operatorname{tg}\alpha<0$. Тому $\sin\alpha\cdot\cos\alpha<0$ (додатне число помножили на від'ємне). Решта виразів додатні: $-\cos\alpha>0$; $\dfrac{\operatorname{tg}\alpha}{\cos\alpha}>0$ (від'ємне поділили на від'ємне); $\sin\alpha-\cos\alpha>0$ (від додатного віднімають від'ємне).\quad{\small\color{gray!80!black}(Знаки тригонометричних функцій)}\par\nopagebreak

\noindent{\footnotesize\color{gray!85!black}Типові помилки: А -- вважали синус тупого кута від'ємним (як косинус); Б -- мінус перед косинусом прийняли за знак значення, хоча $\cos\alpha<0$, тож $-\cos\alpha>0$; В -- частку двох від'ємних чисел ($\operatorname{tg}\alpha<0$, $\cos\alpha<0$) вважали від'ємною; Д -- не врахували, що від додатного числа віднімають від'ємне: $\sin\alpha-\cos\alpha>0$.}\par\vspace{0.15cm}

\par\noindent\textbf{14.} Трапеція рівнобічна, тому $\angle CAD=\angle BDA$ і трикутник $AOD$ рівнобедрений; оскільки $\angle AOD=90^\circ$, то $\angle CAD=45^\circ$. Проведемо висоту $CH$: $AH=\frac{AD+BC}{2}=\frac{14+6}{2}=10$~см, а в прямокутному рівнобедреному трикутнику $ACH$ $CH=AH=10$~см. $S=\frac{6+14}{2}\cdot10=100$~см$^2$.\quad{\small\color{gray!80!black}(Трапеція: площа)}\par\nopagebreak

\noindent{\footnotesize\color{gray!85!black}Типові помилки: А -- вважали кут при основі рівним $45^\circ$: висота $\frac{AD-BC}{2}=4$; Б -- висотою взяли лише висоту трикутника $AOD$: $\frac{AD}{2}=7$; Г -- площу рахували як $AD\cdot h$; Д -- площу трапеції рахували без множника $\frac12$.}\par\vspace{0.15cm}

\par\noindent\textbf{15.} $\frac{x+y}{xy}=\frac34$, тому $xy=8$; з $x+y=6$: $x_0=4$, $y_0=2$, різниця $2$.\quad{\small\color{gray!80!black}(Нелінійна система: сума й добуток ($\tfrac1x+\tfrac1y$ і $x+y$))}\par\nopagebreak

\noindent{\footnotesize\color{gray!85!black}Типові помилки: А -- знайшли лише $x_0$; Б -- вважали $(x-y)^2=(x+y)^2-2xy$; В -- знайшли $xy=8$; Г -- переплутали $x_0$ і $y_0$.}\par\vspace{0.15cm}

\par\noindent\textbf{16.} 1) $\sqrt{4^3}=8$; 2) $\frac{a(a-2)}{a-2}=a=4$; 3) $\log_{0{,}5}4=-2$.\quad{\small\color{gray!80!black}(Значення виразу за заданого значення змінної)}\par\nopagebreak

\noindent{\footnotesize\color{gray!85!black}Типові помилки: 1~--~Г: помножили основу на показник: $4\cdot\frac32$; 2~--~Д: «скоротили» почленно: $\frac{a^2}{a}+\frac{2a}{2}=2a$; 3~--~Б: не врахували, що основа $0{,}5<1$.}\par\vspace{0.15cm}

\par\noindent\textbf{17.} 1) $|{-x}|-3=|x|-3$ -- парна; 2) $(-x)^3=-x^3$ -- непарна; 3) $(x-2)^2=0$ лише при $x=2$.\quad{\small\color{gray!80!black}(Графік функції -- його властивість: перетин з осями, симетрія)}\par\nopagebreak

\noindent{\footnotesize\color{gray!85!black}Типові помилки: 1~--~А: помилилися в знаку: $|0|-3=3$; 2~--~В: сплутали: кубічна парабола симетрична відносно осі $y$; 3~--~Д: помилилися в знаку зсуву: вісь симетрії $x=-2$.}\par\vspace{0.15cm}

\par\noindent\textbf{18.} 1) $\angle BAD=60^\circ+40^\circ=100^\circ$, тож $\angle ABC=180^\circ-100^\circ=80^\circ$. 3) $AK=AB=AD$, отже, трикутник $AKD$ рівнобедрений з основою $KD$: $\angle AKD=\frac{180^\circ-40^\circ}{2}=70^\circ$. 2) $\angle BKD=\angle BKA+\angle AKD=60^\circ+70^\circ=130^\circ$.\quad{\small\color{gray!80!black}(Кути на рисунку: правильні й рівнобедрені трикутники в многокутнику)}\par\nopagebreak

\noindent{\footnotesize\color{gray!85!black}Типові помилки: 1~--~Д: вважали, що кут $A$ ромба дорівнює $40^\circ$: $180^\circ-40^\circ$; 2~--~Б: знайшли лише кут $AKD$, не додавши кут $BKA$ правильного трикутника; 3~--~А: сплутали основу рівнобедреного трикутника $AKD$: вважали $\angle AKD=\angle KAD$; 3~--~Д: не поділили на $2$: $180^\circ-40^\circ$.}\par\vspace{0.15cm}

\par\noindent\textbf{19.} $s=\int_0^4(3t+2)\,dt=\left(\frac{3t^2}{2}+2t\right)\Big|_0^4=24+8=32$~м.\quad{\small\color{gray!80!black}(Інтеграл від многочлена й степенів (зокрема після розкриття дужок))}\par\nopagebreak

\noindent{\footnotesize\color{gray!85!black}Типові помилки: $14$ -- знайшли швидкість у момент $t=4$; $56$ -- помножили швидкість у кінці на час; $24$ -- загубили доданок $2$.}\par\vspace{0.15cm}

\par\noindent\textbf{20.} $15+x+2x+25=100$, тож $x=20$ і біологію обрали $40\,\%$. $36$ учнів -- це $15\,\%$, тому опитали $36:0{,}15=240$ учнів; біологію обрали $0{,}4\cdot240=96$.\quad{\small\color{gray!80!black}(Кругова діаграма: частки з невідомим, кількість за відомою часткою)}\par\nopagebreak

\noindent{\footnotesize\color{gray!85!black}Типові помилки: $48$ -- знайшли кількість тих, хто обрав англійську ($x\,\%$); $40$ -- записали відсоток, а не кількість учнів; $144$ -- вважали, що біологію обрали $100-15-25=60\,\%$; $72$ -- вважали, що біологію обрали вдвічі більше учнів, ніж історію.}\par\vspace{0.15cm}

\par\noindent\textbf{21.} $\frac{\pi r^2H}{2\pi rH}=\frac r2=\frac{400}{160}$, $r=5$, $H=16$. Діагональ квадрата $10$, сторона $5\sqrt2$, площа $50$; висота призми $10$; $V=500$~см$^3$.\quad{\small\color{gray!80!black}(Правильна призма чи куб і тіло обертання)}\par\nopagebreak

\noindent{\footnotesize\color{gray!85!black}Типові помилки: $1000$ -- радіус циліндра взяли за радіус вписаного кола: сторона $2r$; $800$ -- висотою призми взяли висоту циліндра; $250$ -- сторону основи взяли рівною радіусу.}\par\vspace{0.15cm}

\par\noindent\textbf{22.} Корені чисельника $x=a\pm3$ (завжди різні); знаменник $=(x-4)(x+3)$. Корінь єдиний, коли один із $a\pm3$ дорівнює $4$ або $-3$: $a=1,\,7,\,-6,\,0$; обидва вилучитися не можуть. Сума $2$.\quad{\small\color{gray!80!black}(Дріб дорівнює нулю: корінь знаменника вилучає корінь чисельника)}\par\nopagebreak

\noindent{\footnotesize\color{gray!85!black}Типові помилки: $8$ -- урахували лише корінь знаменника $x=4$; $-6$ -- урахували лише корінь знаменника $x=-3$; $-2$ -- розклали знаменник з помилкою знака: $(x+4)(x-3)$.}\par\vspace{0.15cm}

\end{document}
